% OWD, 2023.

\documentclass{beamer}
\usepackage{tikz-cd,adjustbox,listings,caption,mathtools}

\tikzcdset{outer sep=2pt}

\usetheme{CambridgeUS}
\setbeamertemplate{caption}[numbered]
\numberwithin{figure}{section}

\DeclareMathOperator{\catobj}{ob}
\DeclareMathOperator{\catid}{id}
\DeclareMathOperator{\cathom}{hom}

% Portability note: \text{\textbf{...}} is more universally accepted than
% \mathbf{...} when working with Latin fonts, especially in a Sans typeface as
% with Beamer. I'm still using DeclareMathOperator to ensure correct spacing,
% particularly around parentheses.
\DeclareMathOperator{\catendo}{\text{\textbf{Endo}}}
\DeclareMathOperator{\cathask}{\text{\textbf{Hask}}}
\DeclareMathOperator{\catset}{\text{\textbf{Set}}}
\DeclareMathOperator{\catpos}{\text{\textbf{Pos}}}
\DeclareMathOperator{\rvect}{\mathbb{R} \text{\textbf{-Vect}}}

\lstset{language = Haskell,
        showstringspaces = false,%
        columns = flexible,%
        basicstyle = {\small\ttfamily},%
        keywordstyle = \color{blue},%
        stringstyle = \color{darkgray}
}

\title[Elementary Category Theory]%
{An Investigation of Elementary Category Theory}

\subtitle{\ldots\ \fontfamily{lmr}%
        \textsc{with Selected Applications in PM and TCS.}}

\author[Dixon, Drury \& Brook]%
{Oliver~Dixon \and Matthew~Drury \and Ben~Brook}

\institute[]{Department of Mathematics, University of York}
\date{Spring Term, 2023}

\AtBeginSection
{
        \begin{frame}
                \frametitle{Presentation Overview}
                \tableofcontents[currentsection]
        \end{frame}
}

\begin{document}

\titlegraphic{\adjustbox{scale=0.7}{%
        % This could be in the preamble, but most editors that pseudo-parse the
        % TeX to provide syntax highlighting won't consider maths environments
        % unless they're inside the document proper.
        \begin{tikzcd}[sep=large, color=gray, ampersand replacement=\&]
                A
                        \arrow[r, "\alpha"]
                        \arrow[d, "\mu" swap] \&
                A^\prime
                        \arrow[r, "\alpha^\prime"]
                        \arrow[d, "\mu^\prime" swap] \&
                A^{\prime\prime}
                        \arrow[r, "\alpha^{\prime\prime}"]
                        \arrow[d, "\mu^{\prime\prime}" swap] \& {} \\
                B \arrow[r, "\beta" swap] \&
                B^\prime \arrow[r, "\beta^\prime" swap] \&
                B^{\prime\prime} \arrow[r, "\beta^{\prime\prime}" swap] \& {}
        \end{tikzcd}
        \textcolor{gray}{\Large\ldots}%
        \begin{tikzcd}[sep=large, color=gray, ampersand replacement=\&]
                {} \arrow[r, "\alpha^{(n-1)}"] \&
                A^n \arrow[d, "\mu^n"] \\
                {} \arrow[r, "\beta^{(n-1)}" swap] \& B^n
        \end{tikzcd}}%
}

\frame{\titlepage}

\section{Theoretical Underpinnings: Axiomatic Constructions [MD]}

\section{Category-Theoretic Interpretations of Familiar Structures [BB]}
% category of sets
% poset category
% vector spaces category
% category of proofs
\begin{frame}
        \frametitle{Category of Sets}
        Let $A$ be a set and $f \colon A \to A$ be a function. We will denote
        $1_A \coloneqq f$.  Now let $f \colon A \to B$, $y \colon B \to C$, and
        $h \colon C \to D$ with $A$, $B$, $C$, and $D$ being sets. We have that
        \begin{equation*}
                \left( f \circ g \right) \circ h = f \circ \left( g \circ h
                \right) .
        \end{equation*}
        Hence the category axioms are satisfied, and there is a category whose
        objects are all the sets and morphisms are all the functions between
        those sets. We'll denote this category $\catset$ moving forwards.
\end{frame}

\begin{frame}
        \frametitle{Category of Posets}
        A \emph{poset} is a set $A$ with a relation, $\sim _A$, that is
        reflexive, transitive, and antisymmetric. A monotone function between
        posets $A$ and $B$, $m \colon A \to B$, is such that
        \begin{equation*}
                \forall a_1, a_2 \in A, \, a_1 \sim _A a_2 \implies m(a_1) \sim
                _B m(a_2).
        \end{equation*}
        We have a category $\catpos$ whose objects are posets and morphisms are
        monotone functions between those posets.
\end{frame}

\begin{frame}
        \frametitle{Category of Vector Spaces}
        We have a category, $\rvect$, whose objects are \emph{vector spaces} over
        $\mathbb{R}$ and morphisms are $\mathbb{R}$\emph{-linear maps} between
        those vector spaces. We take a linear map to be a function, $f$, between
        two vector spaces that preserves vector addition and scalar
        multiplication. That is,
        \[
                f( \vec{u} + \vec{v} ) = f( \vec{u} ) + f( \vec{v} ) ,
        \]
        and,
        \[
                f( a \vec{u} ) = af( \vec{u} ) .
        \]
\end{frame}

\section[Further Applications: Func. Prog. \protect\& Lambda-Calculus {[OD]}]%
        {Further Applications: Functional Programming and
        \texorpdfstring{$\lambda$}{Lambda}-Calculus [OD]}
\begin{frame}
        \frametitle{Monoidal Categories (1)}
        To construct a monoidal category $\mathcal{C}$, we need to consider six
        elements:
        \begin{itemize}[<+->]
                \item A base category $\mathcal{C}_0$
                \item A bifunctor $\otimes:\mathcal{C}_0 \times \mathcal{C}_0
                        \to \mathcal{C}_0$
                \item An identity from our base category, $I \in \catobj
                \mathcal{C}_0$
                \item An associativity natural transformation:
                        $\alpha_{A,B,C} : (A \otimes B) \otimes C \to A \otimes
                        (B \otimes C)$
                \item A left-identity natural transformation:
                        $\lambda_A: I \otimes A \to A$
                \item A right-identity natural transformation:
                        $\rho_A: A \otimes I \to A$.
        \end{itemize}
        \onslide<+->{Then, $\mathcal{C}=(\mathcal{C}_0, \otimes, I, \alpha,
        \lambda, \rho)$, or just $\mathcal{C}=(\mathcal{C}_0, \otimes, I)$.}
\end{frame}

\begin{frame}
        \frametitle{Monoidal Categories (2)}
        \begin{figure}
                \adjustbox{scale=0.8, center}{%
                \begin{tikzcd}[sep=large, ampersand replacement=\&]
                        \left[(A \otimes B) \otimes C \right] \otimes D
                                \arrow[r, "\alpha"]
                                \arrow[d, "\alpha\,\otimes\,\catid D" swap] \&
                        (A \otimes B) \otimes (C \otimes D)
                                \arrow[r, "\alpha"] \&
                        A \otimes \left[ B \otimes (C \otimes D) \right] \\
                        \left[ A \otimes (B \otimes C) \right] \otimes D
                                \arrow[rr, "\alpha" swap] \& \&
                        A \otimes \left[ (B \otimes C) \otimes D \right]
                                \arrow[u, "\catid A\,\otimes\,\alpha" swap]
                \end{tikzcd}}
                \caption{Associativity induced by $\alpha$}
        \end{figure}

        \vfill
        \pause
        \begin{figure}
                \adjustbox{scale=0.8, center}{%
                \begin{tikzcd}[sep=large, ampersand replacement=\&]
                        (B \otimes I) \otimes C
                                \arrow[rr, "\alpha"]
                                \arrow[rd, "\rho\,\otimes\,\catid C" swap] \& \&
                        B \otimes (I \otimes C)
                                \arrow[ld, "\catid B\,\otimes\,\lambda"] \\
                        \& B \otimes C \&
                \end{tikzcd}}
                \caption{Left- and right-identities induced by $\lambda$ and
                $\rho$}
        \end{figure}
\end{frame}

\begin{frame}
        \frametitle{Monoids (1)}
        A monoid $(M,\mu,\eta)$ is composed of:
        \begin{itemize}[<+->]
                \item Some object $M \in \catobj \mathcal{C}_0$
                \item An associated ``multiplication'' bifunctor $\mu \in
                        \cathom (M \otimes M, M)$
                \item An associated ``unit'' identity $\eta \in \cathom (I, M)$.
        \end{itemize}
\end{frame}

\begin{frame}
        \frametitle{Monoids (2)}

        ($\alpha$, $\lambda$, and $\rho$ refer to the natural isomorphic
        transformations from the parent monoidal category.)

        \begin{figure}
                \adjustbox{scale=0.8, center}{%
                \begin{tikzcd}[sep=large, ampersand replacement=\&]
                        (M \otimes M) \otimes M
                                \arrow[r, "\alpha"]
                                \arrow[d, "\mu\,\otimes\,\catid M" swap] \&
                        M \otimes (M \otimes M)
                                \arrow[r, "\catid M\,\otimes\,\mu"] \&
                        M \otimes M
                                \arrow[d, "\mu"] \\
                        M \otimes M
                                \arrow[rr, "\mu" swap] \& \& M
                \end{tikzcd}}
                \caption{Associativity of the monoid}
        \end{figure}

        \vfill
        \pause
        \begin{figure}
                \adjustbox{scale=0.8, center}{%
                \begin{tikzcd}[sep=large, ampersand replacement=\&]
                        I \otimes M
                                \arrow[r, "\eta\,\otimes\,\catid M"]
                                \arrow[rd, "\lambda" swap] \&
                        M \otimes M
                                \arrow[d, "\mu" swap] \&
                        M \otimes I
                                \arrow[l, "\catid M\,\otimes\,\eta" swap]
                                \arrow[ld, "\rho"] \\
                        \& M \&
                \end{tikzcd}}
                \caption{Left- and right-identities of the monoid}
        \end{figure}
\end{frame}

\begin{frame}
        \frametitle{Functional Programming Languages}
        \begin{itemize}[<+->]
                \item In \textit{purely} functional programming languages,
                        functions must not cause side-effects: for some input,
                        they must always return the same output, independent of
                        the state of the wider environment.
                \item There is no allowance for a shared mutable state. At first
                        glance, this causes problems, since many common
                        operations that are made easy in the imperative world,
                        are in direct contravention with this safety principle.
                \item How can monoids and monodial categories help us with
                        printing ``Hello World''?
        \end{itemize}
\end{frame}

\begin{frame}
        \frametitle{An Especially Useful Monoidal Category}
        Let's consider an example: the category of \textit{endofunctors}! We
        have already encountered endofunctors in the first section.
        \begin{itemize}[<+->]
                \item The endofunctors form a (strict) monoidal category.
                \item The normal functor composition operation becomes the
                        associated bifunctor $\otimes$.
                \item The identity functor becomes the identity element $I$.
                \item By our previous definitions, an endofunctor can be
                        interpreted as a monoid.
                \item A \textit{monad} is a monoid in this category.
                \item We can specify the base category $\mathcal{C}_0$ with
                        $\catendo{(\mathcal{C}_0)}$.
        \end{itemize}
\end{frame}

\begin{frame}
        % I can't use a math operator (\cathask or \catendo) in a title that may
        % be written to a PDF string or index; some viewers will not like it!
        \frametitle{Haskell, \textbf{Hask}, and \textbf{Endo}(\textbf{Hask})}
        \begin{center}
                \begin{minipage}{.2\textwidth}
                        \includegraphics{hasklogo.pdf}
                \end{minipage}
                \begin{minipage}{.7\textwidth}
                        The \textit{Haskell} language is a useful demonstration
                        tool, however the type systems of most languages will be
                        equally valid for these purposes. (Early versions of
                        Haskell weren't capable of I/O!)
                \end{minipage}
        \end{center}
        \pause
        \begin{itemize}[<+->]
                \item Consider the category $\cathask$: its objects are Haskell
                        types, and its morphisms are functions between types.
                        Haskell is a \textit{purely} functional language, so
                        $\cathask$ is well-defined.
                \item What about $\catendo(\cathask)$?
                \item Its objects are endofunctors on $\cathask$, and its
                        morphisms are natural transformations between these
                        functors.
                \item Thus, endofunctors on $\cathask$ are monads.
        \end{itemize}
\end{frame}

\begin{frame}
        \frametitle{In Closing: Practical Notes}
        We have established a tangential, perhaps contrived, link between
        Category Theory and functional programming, the latter of which is 
        based on the $\lambda$-calculus. Why is this interdisciplinary
        observation useful?
        \pause
        \begin{itemize}[<+->]
                \item In practice, monads allow stateful calculation. Unlike
                the mindless emulation of imperative languages, this is achieved
                without breaking function purity. This can be seen intuitively
                as \textit{threading state though a chain of functions}.
                \item This is akin to ``following the arrows'' on a commutative
                diagram of $\catendo{(\cathask)}$ using the category product
                $\cathask \times \cathask$ to pass the desired parameter with
                the current state, and returning the transformed state,
                recalling that the endofunctors are morphisms in $\cathask$.
        \end{itemize}
\end{frame}

\begin{frame}[fragile]
        \frametitle{Any Questions?}

        \begin{figure}
                \adjustbox{scale=0.7, center}{%
                \begin{lstlisting}[frame=single,linewidth=.76\textwidth,%
                                gobble=23]
                        main :: IO () -- IO is a standard monad!
                        main = do
                                putStrLn ( "What is your name?" )
                                name <- getLine
                                if name /= "QUIT" then do
                                        putStrLn ( "Hello, " ++ i ++ "." )
                                else
                                        return ()
                \end{lstlisting}}
                \caption*{``A Haskell User-Greeting Machine''}
        \end{figure}
\end{frame}
\end{document}

